\documentclass[journal,10pt]{IEEEtran}
\setlength{\textwidth}{182mm}
\setlength{\columnsep}{4.2mm}
\usepackage{amsmath,amssymb,newtxmath}
\usepackage[no-math]{fontspec}
\setmainfont{Times New Roman}
\usepackage{xeCJK}
\setCJKmainfont[AutoFakeBold=2]{SimSun}
\usepackage{graphicx,booktabs,array,hyperref,placeins}
\hypersetup{hidelinks,unicode=true}
\raggedbottom
\setlength{\emergencystretch}{2em}
\renewcommand{\arraystretch}{1.15}
\usepackage{seqsplit}
\begin{document}
\title{Litz 线的分层验证、共端复电流分配与有限长度纵向 PEEC · 02 方法}
\author{缪雨轩\thanks{淮南师范学院}}

\maketitle
\noindent 合并稿对应章节 / Chapter in merged manuscript: \href{paper_zh-CN_ieee.pdf#chapter-02}{02}\par
\FloatBarrier
\section*{02 方法}
\hypertarget{chapter-02}{}
\pdfbookmark[0]{02 方法}{bookmark-chapter-02}
\subsection*{工况、物理量与公平比较}
\hypertarget{02.conventions.heading}{}
\pdfbookmark[1]{工况、物理量与公平比较}{bookmark-02.conventions.heading}
计算采用 SI 单位、exp(jωt) 和 RMS 相量。材料为 20 ℃ 铜，σ=5.8×10⁷ S/m、μr=1；主输运工况 f=200 kHz、总电流 20 A RMS，无独立外加场。股线沿途绝缘、端部并联。所有“每米”指轴向长度，实际股线弧长另行报告。有限段端电压 V 与二维模型的单位轴向长度端场具有不同单位，比较时必须保留长度归一化。


原题 Q1 的 D=2 mm 圆线保留为解析校准，其铜面积为 3.141593 mm²、Jdc=6.366198 A/mm²，超过 4 A/mm² 约束。它不是满足最终设计条件的导体。Q2 的名义铜面积至少 5 mm²；640×0.100 mm 的理想圆面积为 5.026548 mm²，而原二维 64 边形离散面积约 5.018478 mm²。采用实际面积的等铜参照，才能避免把几何面积误差解释成纯电磁改善。


不同股长的直流电阻由并联电导相加得到 \href{paper_zh-CN_ieee.pdf#E01}{(1)}。同时报告铜体积、弧长增长、绝对 Rac、Rac/Rdc、损耗和逐股复电流，而不是只给比值。RMS 焦耳损耗采用 \href{paper_zh-CN_ieee.pdf#E02}{(2)}，没有峰值相量的 1/2 因子。存在独立外场时，吸收功率含外场做功；RT0 外场损耗不能无条件除以 20² 命名为内禀输运电阻。几何等铜也不意味着单位轴向长度的铜体积相同。


\subsection*{分层模型与验证方法}
\hypertarget{02.hierarchy.heading}{}
\pdfbookmark[1]{分层模型与验证方法}{bookmark-02.hierarchy.heading}
模型层级按未知量与可比较观测量划分。解析圆线给出长直均质导体的内部阻抗和径向电流；二维扩散验证相同内场问题。二维共端多股模型进一步保留股内分布、股间复电流及明确回流，但省略纵向换位与有限端区。纵向 PEEC 保留有限长度部分电感核和共端约束，却不包含完整横向环流。三分量 RT0 外场算例检验三维算子、连续性和外场吸收，仍与曲线 Litz 的输运问题不同。


验证同时包括求解残差、约束、独立功率与参数敏感性。真实线性残差和公共端/总电流约束沿用 10⁻⁸ 门槛，独立功率平衡为 0.1\%，相关网格、分段、边界和长度变化为 1\%。解析参考及解析电流/功率积分有更严格的专门门槛；RT0 GMRES 目标保持 10⁻¹¹，积分敏感性为 0.1\%；CPU/GPU 同矩阵与复电流对照为 10⁻¹⁰。这些门槛属于本研究的检查定义，不是全部物理问题已经具有同样精度的宣言。


证据可在共同假设下交叉比较，不能沿箭头自动升级结论。小线性残差说明已装配系统被正确求解；它不排除遗漏物理或错误边界。功率平衡约束离散方程与后处理的一致性，但保守的错误模型也可能平衡。两档结果接近限制特定扰动下的变化，不能给出无条件误差上界或统计置信区间。归档因此把执行状态、数值检查状态和物理适用范围分别保存。


\subsection*{共端场路与复电流均流判据}
\hypertarget{02.common.heading}{}
\pdfbookmark[1]{共端场路与复电流均流判据}{bookmark-02.common.heading}
端部并联并不要求各股预先具有相同电流。令 Z 为指定线性模型的耦合阻抗矩阵，e 为独立外场的感应项，i 为股线复电流向量，则公共端电压和总电流形成增广系统 \href{paper_zh-CN_ieee.pdf#E04}{(3)}。实际实现求解增广系统，不依赖显式矩阵求逆。约束允许幅值和相位由场路共同决定；仅比较 |ik| 而忽略相位，会漏掉复电流和总电流之间的重要关系。


在增广系统解唯一的前提下，把 i=(It/N)1 代入端口方程，均流成立当且仅当 [(It/N)Z1+e] 落在 1 的张成空间。用投影 Q 表示即得到 \href{paper_zh-CN_ieee.pdf#E05}{(4)}。无独立外场时，这成为阻抗行和相同的条件；外场存在时，感应项可能保持或破坏该条件。这个代数判据适用于所指定的 Z、e 和端接，不是任意连续三维换位结构的均流定理。


二维共端 FEM 将公共轴向端场与 Az 自由度一起求解。保存的 u 与端场具有 V/m 单位，Jz=σ(v−u)；回流使用自己的端场，空气不赋导体电流。图件从原节点和材料三角形恢复 Jz，按原毫米坐标缩放执行积分。独立线性三角形质量积分与各股复电流、束内损耗、回流损耗和铜面积核对，因此颜色图的来源可审查。保存的 u 不能直接当作电流密度绘制。


\subsection*{有限长度纵向 PEEC：理论与离散}
\hypertarget{02.peec_method.heading}{}
\pdfbookmark[1]{有限长度纵向 PEEC：理论与离散}{bookmark-02.peec_method.heading}
纵向模型以截面体积细丝承载轴向、沿长度均匀的复电流。欧姆项为各细丝的有限段电阻，部分电感由有限体积对的 1/|r−r′| 核积分得到 \href{paper_zh-CN_ieee.pdf#E08}{(5)}，公共端口和总电流仍按 \href{paper_zh-CN_ieee.pdf#E04}{(3)} 求解。这里的 1264 与 5024 是体积细丝自由度，不是相同数量的绝缘 Litz 股。场内径向电流集中可以由截面电流分配表现，但基函数不含完整横向涡流环路或真实端部重分布。


互项积分阶数与自项精度作为独立选项归档。1264 自由度固定截面首先使用 0.16、0.32、0.64 m 长度；积分检查固定长度，在互项 q=4、6、8 和自项 10/8\(\to\)12/10 之间分别改变一个因素。5024 代表模型复核末两档长度和积分。由于 m40/m80 的实际铜面积略有不同，跨 m 对照按每档自身的等面积 Bessel 参照解释，不能把面积变化隐藏为纯离散误差。


部分电感中的外部磁场与长直圆线的内部 Bessel 阻抗不同，因此这里主要以 Rac 作解析比较，并保留长度和几何面积。长度变化限制了当前有限核计算对所选长度的敏感性，却没有增加新的横向基函数。与 FEM–PEEC 文献 [4] 比较时，本研究的纵向体积细丝、激励、回流与反馈方式应分别注明；某个耦合实现的假设不能推广为所有 PEEC 方法的共同缺点。


\hypertarget{E01}{}
\begin{equation}
\begin{gathered}
R_{\mathrm{dc}}=\left(\sum_{k=1}^{N}\frac{\sigma A_k}{\ell_k}\right)^{-1} \\
V_{\mathrm{Cu}}=\sum_{k=1}^{N}A_k\ell_k
\end{gathered}
\tag{1}\label{eq:E01}
\end{equation}
\hypertarget{E02}{}
\begin{equation}
\begin{gathered}
P_{\mathrm{Cu}}=\int_{\Omega_{\mathrm{Cu}}}|\mathbf J|^2/\sigma\,\mathrm d\Omega \\
R_{\mathrm{ac}}=P_{\mathrm{Cu}}/|I_t|^2 \\
r_{\mathrm{ac}}=R_{\mathrm{ac}}/L_z
\end{gathered}
\tag{2}\label{eq:E02}
\end{equation}
\hypertarget{E04}{}\begin{equation}\begin{bmatrix} Z&-\mathbf1\\ \mathbf1^T&0\end{bmatrix}\begin{bmatrix}\mathbf i\\V\end{bmatrix}=\begin{bmatrix}-\mathbf e\\I_t\end{bmatrix}\tag{3}\label{eq:E04}\end{equation}
\hypertarget{E05}{}
\begin{equation}
\begin{gathered}
Q=I_N-\mathbf1\mathbf1^T/N \\
Q\left[(I_t/N)Z\mathbf1+\mathbf e\right]=\mathbf0
\end{gathered}
\tag{4}\label{eq:E05}
\end{equation}
\hypertarget{E08}{}
\begin{equation}
\begin{gathered}
R_{\alpha\beta}=\delta_{\alpha\beta}\frac{L_z}{\sigma A_\alpha} \\
L_{\alpha\beta}=\frac{\mu}{4\pi A_\alpha A_\beta}\int_{V_\alpha}\int_{V_\beta}\frac{\hat{\mathbf t}_\alpha\cdot\hat{\mathbf t}_\beta}{|\mathbf r-\mathbf r'|}\,\mathrm dV\,\mathrm dV'
\end{gathered}
\tag{5}\label{eq:E08}
\end{equation}
\begin{figure}[!htbp]
\centering
\setcounter{figure}{0}
\refstepcounter{figure}\label{fig:model_hierarchy}
\hypertarget{fig-model_hierarchy}{}
\includegraphics[width=88.9mm]{../../../panels/zh-CN/model_hierarchy.pdf}
\par\vspace{4pt}{\small\centering 图 1. 分层模型与可比较观测量。\par}
\vspace{3pt}{\small\raggedright 输运模型按物理对象与观测量分组：圆线内场用于解析与二维校准，多股二维场路求解共同端口下的复电流，纵向 PEEC 检验有限长度效应。连线标明内部场校准或等面积电阻参照。RT0 外场损耗与三维 A–φ 场诊断各自构成独立研究范围。\par}
\end{figure}
\begin{figure}[!htbp]
\centering
\setcounter{figure}{1}
\refstepcounter{figure}\label{fig:common_terminal}
\hypertarget{fig-common_terminal}{}
\includegraphics[width=88.9mm]{../../../panels/zh-CN/common_terminal.pdf}
\par\vspace{4pt}{\small\centering 图 2. 多股共同端口与回流。\par}
\vspace{3pt}{\small\raggedright 各股沿途绝缘，在两端并联，逐股复电流由共同端电压及总电流约束求解。耦合矩阵 Z 与激励项 e 包含指定模型中的阻抗和场作用。受控回流承担反向总电流，并具有独立的端场。图中给出已知输入及待求复电流和电压。\par}
\end{figure}
\clearpage
\begin{thebibliography}{99}
\bibitem[4]{P03} Andreas Roßkopf, Eberhard Bär, Christopher Joffe, Clemens Bonse. Calculation of Power Losses in Litz Wire Systems by Coupling FEM and PEEC Method. IEEE TPEL 31(9):6442–6449, 2016. https://doi.org/10.1109/TPEL.2015.2499793
\end{thebibliography}
\end{document}
